\chapter{1988:Johann Deisenhofer and Robert Huber and Hartmut Michel}

\label{chap:chemistry1988}

The Nobel Prize in Chemistry for the year 1988 was awarded jointly to Johann Deisenhofer and Robert Huber and Hartmut Michel.

\textbf{Motivation:}

for the determination of the three-dimensional structure of a photosynthetic reaction centre

\section{Johann Deisenhofer}

\subsection*{Early Life and Education}

Johann Deisenhofer was born in 1943 in Zusamaltheim, Germany.

\subsection*{Career and Major Research}

Deisenhofer studied physics at the Technical University of Munich and earned his Ph.D. in 1974 at the Max Planck Institute of Biochemistry in Martinsried under Robert Huber. He later joined the University of Texas Southwestern Medical Center in Dallas.

\subsection*{Nobel Prize-winning Work}

The work for which Johann Deisenhofer was awarded the Nobel Prize in Chemistry in 1988 was described by the Nobel Committee as: "for the determination of the three-dimensional structure of a photosynthetic reaction centre".

He determined the three-dimensional structure of the photosynthetic reaction centre of the purple bacterium Rhodopseudomonas viridis by X-ray crystallography, showing how light is captured and converted into a charge separation across the membrane.

\subsection*{Significance and Impact}

This was the first structure of an integral membrane protein determined at atomic resolution, and it revealed in detail how photosynthesis works at the molecular level.

\subsection*{Later Career and Honors}

Deisenhofer continued his research as an investigator of the Howard Hughes Medical Institute at the University of Texas Southwestern Medical Center. He shared the 1988 Nobel Prize in Chemistry.

\subsection*{Legacy}

The reaction-centre structure became the reference model for membrane protein crystallography.

\subsection*{Selected Publications}

"Structure of the Protein Subunits in the Photosynthetic Reaction Centre of Rhodopseudomonas viridis at 3 Å Resolution" (Nature, 1985)

\clearpage

\section{Robert Huber}

\subsection*{Early Life and Education}

Robert Huber was born in 1937 in Munich, Germany.

\subsection*{Career and Major Research}

Huber studied chemistry at the Technical University of Munich, where he received his Ph.D. in 1963. He became a director at the Max Planck Institute of Biochemistry in Martinsried.

\subsection*{Nobel Prize-winning Work}

The work for which Robert Huber was awarded the Nobel Prize in Chemistry in 1988 was described by the Nobel Committee as: "for the determination of the three-dimensional structure of a photosynthetic reaction centre".

He led the crystallographic work, using the crystals prepared by Hartmut Michel, that solved the structure of the photosynthetic reaction centre.

\subsection*{Significance and Impact}

The structure determined how the protein complex converts absorbed light into separated charges, and it established X-ray crystallography as a method for studying membrane proteins.

\subsection*{Later Career and Honors}

Huber remained at the Max Planck Institute of Biochemistry in Martinsried. He shared the 1988 Nobel Prize in Chemistry.

\subsection*{Legacy}

His structural studies extended across many proteins and made Martinsried a center of structural biology.

\subsection*{Selected Publications}

"Structure of the Protein Subunits in the Photosynthetic Reaction Centre of Rhodopseudomonas viridis at 3 Å Resolution" (Nature, 1985)

\clearpage

\section{Hartmut Michel}

\subsection*{Early Life and Education}

Hartmut Michel was born in 1948 in Ludwigsburg, Germany.

\subsection*{Career and Major Research}

Michel studied at the University of Tübingen and received his Ph.D. from the University of Würzburg in 1977. He conducted research at the Max Planck Institute of Biochemistry in Martinsried and, from 1987, became a director at the Max Planck Institute of Biophysics in Frankfurt.

\subsection*{Nobel Prize-winning Work}

The work for which Hartmut Michel was awarded the Nobel Prize in Chemistry in 1988 was described by the Nobel Committee as: "for the determination of the three-dimensional structure of a photosynthetic reaction centre".

He succeeded in crystallizing the photosynthetic reaction centre of Rhodopseudomonas viridis, a membrane protein that had resisted crystallization, which made the structure determination possible.

\subsection*{Significance and Impact}

Michel's crystallization of the reaction centre was a breakthrough for structural biology and pioneered methods for crystallizing membrane proteins.

\subsection*{Later Career and Honors}

Michel continued his research at the Max Planck Institute of Biophysics in Frankfurt. He shared the 1988 Nobel Prize in Chemistry.

\subsection*{Legacy}

His approach to membrane protein crystallization opened the way to the atomic structures of many other membrane proteins.

\subsection*{Selected Publications}

"Structure of the Protein Subunits in the Photosynthetic Reaction Centre of Rhodopseudomonas viridis at 3 Å Resolution" (Nature, 1985)

\clearpage

\section*{Chapter Summary}

The 1988 Nobel Prize in Chemistry recognized the first atomic structure of a photosynthetic reaction centre. Michel prepared crystals of the membrane protein complex, and Huber and Deisenhofer used them to reveal how light energy is converted into a separated charge across the membrane.
